
Norint apibėžti grupes $\spin$ ir $\pinpm$, prima reikia pasakyti kas yra Clifordo algebros
$\mathup{Cl}^+(3)$ ir $\mathup{Cl}^-(3)$. Šios algebros gaunamos iš vektorinės erdvės $\mathup V = \set{R}^3$.
Tegul $\mathup{T(V)}$ būna erdės $\mathup V$ vienetinė tenzorinė algebra. Jos elementai yra vektorių $v_i \in \mathup V$ tenzorinės
sandaugos $v_1 \otimes v_2 \otimes \dots \otimes v_n$. $\mathup{T(V)}$ algebros idealas $\mathup I^+$ yra generuojamas
sąryšio
\begin{equation}
	\vb v \otimes \vb w
	+
	\vb w \otimes \vb v
	=
	2(\vb v , \vb w) \mathbb 1
	\quad ,
	\label{eq:cl+}
\end{equation}
čia $(\cdot,\cdot)$ yra įprasta $\set R ^3$ skaliarinė sandauga ir $\mathbb 1$ yra algebros $\mathup{T(V)}$
vienetinis elementas, kurį neformaliai galimą suprasti kaip "nulio vektorių tezorinę sandaugą".
Jei pažymime $S \subset \mathup{T(V)}$
\begin{equation}
	S = \qty{
		\vb v \otimes \vb w
		+
		\vb w \otimes \vb v
		-
		2(\vb v , \vb w) \mathbb 1
		\mid
		\vb v, \vb w \in \mathup{T(V)}
	}
\end{equation}
\begin{equation}
	I^+ = \qty{
		\sum_{\substack{
			s \in S \\ 
			a,b \in \mathup{T(V)} 
		}}
		a \otimes s \otimes b
	}
	\quad .
\end{equation}
Tada Clifordo algebra $\mathup{Cl^+(V)}$ apibėžiame kaip faktorimę algebrą $\mathup{T(V)/I^+}$.
Jos elementai yra šoninės klasės $a + \mathup I^+ = {a + i \mid i \in \mathup I^+}$, čia $a$
yra šoninę klasę atstovaujantis elementas iš $\mathup{T(V)}$. Clifordo algebra $\mathup{Cl^-(V)}$
apibrėžiama analogiškai, tik $\mathup I^+$ pakeičiamas idealą $\mathup I^-$, generuojamą iš sąryšio
\begin{equation}
	\vb v \otimes \vb w
	+
	\vb w \otimes \vb v
	=
	- 2(\vb v , \vb w) \mathbb 1
	\quad .
	\label{eq:cl-}
\end{equation}


Kaip galima suprasti tokią konstrukciją? Algebra $\mathup{Cl^\pm(V)}$ yra tenzorinė algebra $\mathup{T(V)}$
kuri "užmiršo" kaip atskirti $\vb v \vb w + \vb w \vb v$ ir $\pm2(\vb v, \vb w) \mathbb 1$,
čia ir toliau nebenaudojame tenzorinės sandaugos simbolio $\otimes$, kad atskirti sandaugą Clifordo algebroje
nuo bendros tenzorinės sandaugos.
Šis užmiršimas leidžia savybes (\ref{eq:cl+}) ir (\ref{eq:cl-}) naudoti supaprastinant išraiškas kaip
\begin{equation}
\vb v \vb v 
= \frac{1}{2}(\vb v \vb v + \vb v \vb v) 
= \pm(\vb v,\vb v) \mathbb 1
\end{equation}
ir
\begin{equation}
\vb v \vb w \vb v 
= \qty( \pm2(\vb v,\vb w)\mathbb 1 - \vb{wv} )\vb v 
= \pm2\vb{(v,w)} \vb v - \vb{wvv} 
=  \pm2\vb{(v,w)} \vb v \mp \vb{(v,v)}\vb w 
\quad .
\end{equation}

Tegul $\qty{\vb e_1,\vb e_2,\vb e_3}$ yra ortonormali $\set R^3$ bazė. Bazinius vektorius galime perketi
į algebra $\mathup{Cl^\pm(V)}$, kur jie turi tenkinti
\begin{equation}
	\begin{split}
		\vb e_i \vb e_j &= - \vb e_j \vb e_i \ , \ i \neq j \\
		\vb e_i \vb e_i &= \pm \mathbb 1
		\quad .
	\end{split}
\end{equation}
Su šių sąryšių pagalba galime bet kokį $\mathup{Cl^\pm(V)}$ elementą užrašyti, kaip
tiesinę kombinaciją aštuonių elementų:
\begin{equation}
		\mathbb 1 ,\,
		\vb e_1 ,\, \vb e_2 ,\, \vb e_3 ,\, 
		\vb e_2 \vb e_3 ,\, \vb e_1 \vb e_3 ,\, \vb e_1 \vb e_2 ,\,
		\vb e_1 \vb e_2 \vb e_3 ,
\end{equation}
nes su pirma taisykle betkurį tenzorinės algebros bazės elementą galima išrikiuoti taip kad
ortų indeksai būtų surašyti didėjimo tvarka ir su antra taisykle pasikartojantys ortai panaikinami.
Tai reiškia $\mathup{Cl^\pm(V)}$ yra aštuonmatės algebros. Algebros elemento norma $\abs{v}$ yra
abibrėžiama, taip pat kaip vektoriaus ilgis
\begin{equation}
	\abs{v} = \sqrt{ \sum_{I \subset {1,2,3}} {v_I}^2 }
	\quad ,
\end{equation}
čia summa eina per visus indeksų poaibius, kurie atitinka bazinius algebgros elementus.
Bendresniu atvėju $\mathup{dim}\qty[\mathup{Cl}^\pm\qty(\set R^n)] = 2^n$.

Apibrėžiame "transponuotą elementą" apibėždami kaip transponavimas veikia ant algebros bazės
\begin{equation}
	\qty(\vb e_{i_1} \vb e_{i_2} \dots \vb e_{i_k} )^t
	=
	\vb e_{i_k} \vb e_{i_{k-1}} \dots \vb e_{i_1}
	=
	(-1)^{k-1} \vb e_{i_1} \vb e_{i_2} \dots \vb e_{i_k} 
	\quad .
\end{equation}
Transpozicija yra algebros anti-homomorfizmas, tai reiškia, kad betkuriems algebros elementams $v$ ir $w$ galioja
\begin{equation}
	(v w)^t = w^t v^t
	\quad .
\end{equation}


Taip pat veikimu ant bazės apibrėžiame algebros homomorfizmą
\begin{equation}
	\alpha( \vb e_{i_1} \vb e_{i_2} \dots \vb e_{i_k} )
	= 
	(-1)^{k} \vb e_{i_1} \vb e_{i_2} \dots \vb e_{i_k} 
	\quad ,
\end{equation}
jis algebros elemento komponentes prie lyginio ortų skaičiau siunčia į savo minusą, o
komponentčių prie nelyginio skaičiaus nepakeičia.











